\chapter{Fourier Transform and rigidity}
\label{chap:3}
\setcounter{section}{-1}

\newcommand{\chthreethmthreezerothree}{\hyperref[thm:3.0.3]{Theorem~3.0.3}}
\newcommand{\chthreethmthreezerofour}{\hyperref[thm:3.0.4]{Theorem~3.0.4}}

\section{Fourier Transform and index of rigidity}\label{sec:3.0}

\sourcepara{3.0.1}
Let \(\ell\) be a prime number. On \(\mathbb{A}^1\) over an algebraically
closed field of characteristic \(\ne \ell\), denote by
\(k:\mathbb{A}^1\to\mathbb{P}^1\) the inclusion. Let
\(K\) in \(\mathrm{D}^b_c(\mathbb{A}^1,\overline{\mathbb{Q}}_\ell)\) be a
perverse sheaf which is a middle extension, i.e., \(K\) is
\(j_*\mathcal{F}[1]\) for a nonempty open set
\(j:U\to\mathbb{A}^1\), and a lisse
\(\overline{\mathbb{Q}}_\ell\)-sheaf \(\mathcal{F}\) on \(U\). We define the
index of rigidity of \(K\), \(\operatorname{rig}(K)\), to be the integer
(compare \hyperref[par:2.0.3]{(2.0.3)}--\hyperref[par:2.0.6]{(2.0.6)})
\[
  \operatorname{rig}(K)
  := \chi\bigl(\mathbb{P}^1,k_*j_*\operatorname{End}(\mathcal{F})\bigr).
\]

\sourceheading[thm:3.0.2]{Theorem 3.0.2}
Let \(p\) and \(\ell\) be prime numbers, \(\ell\ne p\). On
\(\mathbb{A}^1\) over an algebraically closed field of characteristic \(p\),
let \(K\) in
\(\mathrm{D}^b_c(\mathbb{A}^1,\overline{\mathbb{Q}}_\ell)\) be a perverse
irreducible such that neither \(K\) nor \(\operatorname{FT}(K)\) is punctual
(i.e., \(K\) is neither a \(\delta_\alpha\) nor an
\(\mathcal{L}_{\psi_\alpha}[1]\)). Then \(K\) and
\(\operatorname{FT}(K)\) have the same index of rigidity:
\[
  \operatorname{rig}(K)=\operatorname{rig}(\operatorname{FT}(K)).
\]

This is a special case of the slightly more general

\sourceheading[thm:3.0.3]{Theorem 3.0.3}
Let \(p\) and \(\ell\) be prime numbers, \(\ell\ne p\). On
\(\mathbb{A}^1\) over an algebraically closed field of characteristic \(p\),
let \(K\) in
\(\mathrm{D}^b_c(\mathbb{A}^1,\overline{\mathbb{Q}}_\ell)\) be a perverse
sheaf on \(\mathbb{A}^1\) such that both \(K\) and
\(\operatorname{FT}(K)\) are middle extensions from any nonempty open set of
\(\mathbb{A}^1\). Then
\[
  \operatorname{rig}(K)=\operatorname{rig}(\operatorname{FT}(K)).
\]

\begin{proof}[Proof of \chthreethmthreezerothree]
Write \(K=j_*\mathcal{F}[1]\) and
\(\operatorname{FT}(K)=j_*\mathcal{G}[1]\), for some dense open
\(j:U\to\mathbb{A}^1\), and lisse sheaves \(\mathcal{F}\) and
\(\mathcal{G}\) on \(U\). Denote by
\(k:\mathbb{A}^1\to\mathbb{P}^1\) the inclusion. By definition,
\[
\begin{aligned}
  \operatorname{rig}(K)
  &:= \chi\bigl(\mathbb{P}^1,k_*j_*\operatorname{End}(\mathcal{F})\bigr) \\
  &= \dim \operatorname{End}(\mathcal{F})^{I(\infty)}
     + \chi\bigl(\mathbb{A}^1,j_*\operatorname{End}(\mathcal{F})\bigr).
\end{aligned}
\]
The sheaf \(j_*\operatorname{End}(\mathcal{F})\) on \(\mathbb{A}^1\)
contains the subsheaf \(j_*\mathcal{F}\otimes j_*(\mathcal{F}^{\vee})\);
the quotient, supported in \(\mathbb{A}^1-U\), is punctual:
\[
  j_*\operatorname{End}(\mathcal{F})
  \big/
  \bigl(j_*\mathcal{F}\otimes j_*\mathcal{F}^{\vee}\bigr)
  =
  \bigoplus_{x\in\mathbb{A}^1}
  \delta_x\otimes
  \frac{\operatorname{End}(\mathcal{F})^{I(x)}}
       {\mathcal{F}^{I(x)}\otimes(\mathcal{F}^{\vee})^{I(x)}}.
\]
So
\[
\begin{aligned}
  \operatorname{rig}(K)
  &:= \chi\bigl(\mathbb{A}^1,j_*\mathcal{F}\otimes
       j_*(\mathcal{F}^{\vee})\bigr)
     + \dim\bigl(\operatorname{End}(\mathcal{F})^{I(\infty)}\bigr) \\
  &\quad
     + \sum_{x\in\mathbb{A}^1}
       \dim\left[
       \frac{\operatorname{End}(\mathcal{F})^{I(x)}}
            {\mathcal{F}^{I(x)}\otimes(\mathcal{F}^{\vee})^{I(x)}}\right].
\end{aligned}
\]
Since \(K=j_*\mathcal{F}[1]\), \(\mathbb{D}K=j_*\mathcal{F}^{\vee}[1]\),
and we have
\[
  j_*\mathcal{F}\otimes j_*\mathcal{F}^{\vee}
  =
  \mathbb{D}K\otimes K[-2].
\]
Thus
\[
\begin{aligned}
  \operatorname{rig}(K)
  &:= \chi\bigl(\mathbb{A}^1,\mathbb{D}K\otimes K\bigr)
     + \dim\bigl(\operatorname{End}(\mathcal{F})^{I(\infty)}\bigr) \\
  &\quad
     + \sum_{x\in\mathbb{A}^1}
       \dim\left[
       \frac{\operatorname{End}(\mathcal{F})^{I(x)}}
            {\mathcal{F}^{I(x)}\otimes(\mathcal{F}^{\vee})^{I(x)}}\right].
\end{aligned}
\]
Now we break up the \(I(\infty)\)-representation \(\mathcal{F}(\infty)\)
of \(\mathcal{F}\) according to whether slopes are \(\leq 1\) or \(>1\):
\[
  \mathcal{F}(\infty)
  =
  \mathcal{F}(\infty)(\leq 1)\oplus\mathcal{F}(\infty)(>1).
\]
Because there are no nonzero \(I(\infty)\)-equivariant maps between
representations whose sets of slopes are disjoint, we have
\[
\begin{aligned}
  \operatorname{End}(\mathcal{F})^{I(\infty)}
  &:= \operatorname{End}_{I(\infty)}(\mathcal{F}(\infty)) \\
  &=
  \operatorname{End}_{I(\infty)}(\mathcal{F}(\infty)(\leq 1))
  +
  \operatorname{End}_{I(\infty)}(\mathcal{F}(\infty)(>1)).
\end{aligned}
\]
Using this, rewrite the formula for \(\operatorname{rig}(K)\) as a sum of
four terms
\[
\begin{aligned}
  \operatorname{rig}(K)&=1(K)+2(K)+3(K)+4(K),\\
  1(K)&:=\chi\bigl(\mathbb{A}^1,\mathbb{D}K\otimes K\bigr),\\
  2(K)&:=\dim\operatorname{End}_{I(\infty)}
        \bigl(\mathcal{F}(\infty)(>1)\bigr),\\
  3(K)&:=\dim\operatorname{End}_{I(\infty)}
        \bigl(\mathcal{F}(\infty)(\leq 1)\bigr),\\
  4(K)&:=
  \sum_{x\in\mathbb{A}^1}
       \dim\left[
       \frac{\operatorname{End}(\mathcal{F})^{I(x)}}
            {\mathcal{F}^{I(x)}\otimes(\mathcal{F}^{\vee})^{I(x)}}\right].
\end{aligned}
\]
With this breakup, the equality
\(\operatorname{rig}(K)=\operatorname{rig}(\operatorname{FT}(K))\) results
from

\sourceheading[thm:3.0.4]{Theorem 3.0.4}
In the situation of \hyperref[thm:3.0.3]{Theorem~3.0.3}, we have
\[
\begin{aligned}
  1(K)&=1(\operatorname{FT}(K)),&
  2(K)&=2(\operatorname{FT}(K)),\\
  3(K)&=4(\operatorname{FT}(K)),&
  4(K)&=3(\operatorname{FT}(K)).
\end{aligned}
\]
\end{proof}

\begin{proof}[Proof of \chthreethmthreezerofour]
We begin with \((1)\), a form of Parseval's formula valid for any object
\(K\) in \(\mathrm{D}^b_c\). We denote by \(\mathbb{D}_{-}K\) the object
\([x\mapsto -x]^*(\mathbb{D}K)\). We have
\[
\begin{aligned}
  1(K)
  &:=\chi\bigl(\mathbb{A}^1,\mathbb{D}K\otimes K\bigr)
   = \chi_c\bigl(\mathbb{A}^1,\mathbb{D}K\otimes K\bigr) \\
  &=\operatorname{rank}_0\bigl(\mathbb{D}_{-}K*_{!+}K\bigr),
\end{aligned}
\]
the last equality by using proper base change to calculate the stalk at zero
of \(\mathbb{D}_{-}K*_{!+}K\). Applying this to \(\operatorname{FT}K\) gives
\[
\begin{aligned}
  1(\operatorname{FT}K)
  &= \operatorname{rank}_0
     \bigl(\mathbb{D}_{-}\operatorname{FT}K*_{!+}\operatorname{FT}K\bigr) \\
  &= \operatorname{rank}_0
     \bigl(\operatorname{FT}(\mathbb{D}K)*_{!+}\operatorname{FT}K\bigr) \\
  &= \operatorname{rank}_0
     \bigl(\operatorname{FT}(\mathbb{D}K\otimes K)[-1]\bigr) \\
  &= \chi\bigl(\mathbb{A}^1,\mathbb{D}K\otimes K\bigr)
   := 1(K),
\end{aligned}
\]
using \(\mathbb{D}_{-}\circ\operatorname{FT}=\operatorname{FT}\circ
\mathbb{D}\), that \(\operatorname{FT}\) interchanges \(\otimes\) and
\(*_{!+}\), the \(!\)-form of \(\operatorname{FT}\), and proper base change.

Assertion \((2)\) follows from taking dimensions in the more precise equality
\[
  \operatorname{End}_{I(\infty)}(\mathcal{F}(\infty)(>1))
  =
  \operatorname{End}_{I(\infty)}(\mathcal{G}(\infty)(>1)).
\]
This holds because, in terms of Laumon's local Fourier Transform
\(\operatorname{FTloc}(\infty,\infty)\), we have
\[
  \mathcal{G}(\infty)(>1)
  =
  \operatorname{FTloc}(\infty,\infty)(\mathcal{F}(\infty)(>1)),
\]
and one knows (cf. \cite{Lau-TF}, \cite{Ka-TL},
\cite[7.4.1]{Ka-ESDE}) that \(\operatorname{FTloc}(\infty,\infty)\) is an
autoequivalence of the category of \(I(\infty)\)-representations with all
slopes \(>1\).

The last two assertions,
\[
  3(K)=4(\operatorname{FT}(K)),\qquad
  4(K)=3(\operatorname{FT}(K)),
\]
are in fact equivalent. (They are obtained from each other by replacing \(K\)
by \(\operatorname{FT}K\), and noting that
\(\operatorname{FT}(\operatorname{FT}K)\cong [x\mapsto -x]^*K\) up to a Tate
twist, while visibly \(4(K)=4([x\mapsto -x]^*K)\).) We will prove that
\[
  4(K)=3(\operatorname{FT}(K)).
\]
For this, we recall the theory of Laumon's local Fourier Transforms
\(\operatorname{FTloc}(x,\infty)\), for \(x\) in \(\mathbb{A}^1\). Each of
these is an equivalence of categories
\[
  I(x)\text{-representations}
  \cong
  I(\infty)\text{-representations with all slopes }<1.
\]
By Laumon's theory of stationary phase \cite[7.4]{Ka-ESDE}, we know that
\[
  \mathcal{G}(\infty)(\leq 1)
  =
  \bigoplus_{x\in\mathbb{A}^1}
  \mathcal{L}_{\psi_x}\otimes
  \operatorname{FTloc}(x,\infty)
  \bigl(\mathcal{F}(x)/\mathcal{F}(x)^{I(x)}\bigr).
\]
(In \cite[7.4.1]{Ka-ESDE}, the functor
\(M\mapsto
\mathcal{L}_{\psi_x}\otimes\operatorname{FTloc}(x,\infty)(M)\) was called
\(\operatorname{FTloc}(x,\infty)\).)
There are no nonzero \(I(\infty)\)-equivariant maps between the distinct
summands. (If \(M\) and \(N\) are any two \(I(\infty)\)-representations all
of whose slopes are \(<1\), then for \(x\ne y\),
\[
\begin{aligned}
  \operatorname{Hom}_{I(\infty)}
  (\mathcal{L}_{\psi_x}\otimes M,\mathcal{L}_{\psi_y}\otimes N)
  &=
  \operatorname{Hom}_{I(\infty)}
  (M,\mathcal{L}_{\psi_{y-x}}\otimes N)\\
  &=0,
\end{aligned}
\]
since \(\mathcal{L}_{\psi_{y-x}}\otimes N\) has all slopes \(=1\), while
\(M\) has all slopes \(<1\).) Therefore
\[
  \operatorname{End}_{I(\infty)}(\mathcal{G}(\infty)(\leq 1))
  =
  \bigoplus_{x\in\mathbb{A}^1}
  \operatorname{End}_{I(\infty)}
  \left(
  \operatorname{FTloc}(x,\infty)
  \bigl(\mathcal{F}(x)/\mathcal{F}(x)^{I(x)}\bigr)
  \right).
\]
Because \(\operatorname{FTloc}(x,\infty)\) is an equivalence of categories
\[
  I(x)\text{-representations}
  \cong
  I(\infty)\text{-representations with all slopes }<1,
\]
we have
\[
  \operatorname{End}_{I(\infty)}
  \left(
  \operatorname{FTloc}(x,\infty)
  \bigl(\mathcal{F}(x)/\mathcal{F}(x)^{I(x)}\bigr)
  \right)
  =
  \operatorname{End}_{I(x)}
  \bigl(\mathcal{F}(x)/\mathcal{F}(x)^{I(x)}\bigr).
\]
Thus we obtain
\[
  \operatorname{End}_{I(\infty)}(\mathcal{G}(\infty)(\leq 1))
  =
  \bigoplus_{x\in\mathbb{A}^1}
  \operatorname{End}_{I(x)}
  \bigl(\mathcal{F}(x)/\mathcal{F}(x)^{I(x)}\bigr).
\]
To complete the proof, it remains only to check that
\[
  \dim\left[
  \frac{\operatorname{End}_{I(x)}(\mathcal{F}(x))}
       {\mathcal{F}^{I(x)}\otimes(\mathcal{F}^{\vee})^{I(x)}}\right]
  =
  \dim\operatorname{End}_{I(x)}
  \bigl(\mathcal{F}(x)/\mathcal{F}(x)^{I(x)}\bigr).
\]
This is proven in the following section,
\hyperref[prop:3.1.8]{Proposition~3.1.8}.
\end{proof}

\section{Lemmas on representations of inertia groups}\label{sec:3.1}

\sourcepara{3.1.1}
Throughout this section, we fix a complete discrete valuation ring \(R\) with
algebraically closed residue field \(k\) and with fraction field \(K\). We
denote by \(I\) the Galois group
\[
  I:=\operatorname{Gal}(K^{\mathrm{sep}}/K).
\]
If \(\operatorname{char}(k)=p>0\), we denote by \(P\subset I\) the unique
\(p\)-Sylow subgroup. If \(\operatorname{char}(k)=0\), we define
\(P=\{e\}\). The quotient \(I/P\) is (noncanonically) the pro-cyclic group
\(\prod_{\ell\ne p}\mathbb{Z}_{\ell}\).

\sourcepara{3.1.2}
We also fix a prime number \(\ell\ne\operatorname{char}(k)\). By an
\(\ell\)-adic representation of \(I\), we mean a continuous
\(\overline{\mathbb{Q}}_\ell\)-representation \(\rho\) of \(I\) on a
finite-dimensional \(\overline{\mathbb{Q}}_\ell\)-vector space \(V\), with
the property that there exists a finite extension \(E_\lambda\) of
\(\mathbb{Q}_\ell\) with ring of integers \(\mathcal{O}_\lambda\), a free
\(\mathcal{O}_\lambda\)-module \(V_0\), and a continuous
\(\mathcal{O}_\lambda\)-representation \(\rho_0\) of \(I\) on \(V_0\) such
that
\[
  (\rho_0,V_0)\otimes\overline{\mathbb{Q}}_\ell\cong(\rho,V).
\]

\sourcepara{3.1.3}
An \(\ell\)-adic representation of \(I\) is said to be tame if it is trivial
on \(P\). It is said to be unipotent if it is a successive extension of
trivial representations of \(I\). Because the upper unipotent subgroup of
\(\operatorname{GL}(n,\mathcal{O}_\lambda)\) is pro-\(\ell\), any unipotent
representation of \(I\) is tame. If we fix a topological generator \(\gamma\)
of \(I/P\), then isomorphism classes of \(n\)-dimensional unipotent
representations of \(I\) are in bijective correspondence with conjugacy
classes of unipotent elements in
\(\operatorname{GL}(n,\overline{\mathbb{Q}}_\ell)\). By the theory of Jordan
normal form, there is, up to isomorphism, a unique indecomposable unipotent
representation of \(I\) of each dimension \(n\geq 1\), which we denote
\(\operatorname{Unip}(n)\). We sometimes refer to \(\operatorname{Unip}(n)\)
as the ``standard Jordan block of size \(n\)''.

\sourceheading[lem:3.1.4]{Lemma 3.1.4}
If \(M\) and \(N\) are inequivalent irreducible \(\ell\)-adic representations
of \(I\), then \(\operatorname{Ext}^i_I(M,N)=0\) for all \(i\). In particular,
\[
  \operatorname{Hom}_I(M,N)=\operatorname{Ext}^1_I(M,N)=0.
\]

\begin{proof}
\[
\begin{aligned}
  \operatorname{Ext}^i_I(M,N)
  &= \mathrm{H}^i(I,M^{\vee}\otimes N)\\
  &= \mathrm{H}^i(I/P,(M^{\vee}\otimes N)^P).
\end{aligned}
\]
In terms of a topological generator \(\gamma\) of \(I/P\), we have
\[
\mathrm{H}^i(I/P,(M^{\vee}\otimes N)^P)
=
\begin{cases}
  \operatorname{Ker}\bigl(1-\gamma\mid(M^{\vee}\otimes N)^P\bigr),
    & i=0,\\
  \operatorname{Coker}\bigl(1-\gamma\mid(M^{\vee}\otimes N)^P\bigr),
    & i=1,\\
  0, & i\geq 2.
\end{cases}
\]
This explicit description shows that
\(\chi(I,M^{\vee}\otimes N)=0\). Therefore we have
\(\mathrm{h}^0(I,M^{\vee}\otimes N)=\mathrm{h}^1(I,M^{\vee}\otimes N)\),
i.e.,
\[
  \dim\operatorname{Ext}^1_I(M,N)=\dim\operatorname{Hom}_I(M,N).
\]
Since \(M\) and \(N\) are irreducible and inequivalent,
\(\operatorname{Hom}_I(M,N)=0\).
\end{proof}

\sourceheading[lem:3.1.5]{Lemma 3.1.5}
If \(N\) is an irreducible nontrivial \(\ell\)-adic representation of \(I\),
then \(\mathrm{H}^i(I,N)=0\) for all \(i\).

\begin{proof}
Take \(M\) to be the trivial representation in
\hyperref[lem:3.1.4]{Lemma~3.1.4}.
\end{proof}

\sourceheading[lem:3.1.6]{Lemma 3.1.6}
Any \(\ell\)-adic representation \(M\) of \(I\) has a canonical direct sum
``isotypical'' decomposition, indexed by the equivalence classes \(\mathcal{A}\)
of irreducible representations of \(I\), as
\(\bigoplus_{\alpha\in\mathcal{A}} M_\alpha\), where \(M_\alpha\) is a
successive extension of irreducible representations all of type \(\alpha\).

\begin{proof}
This is immediate from \hyperref[lem:3.1.4]{Lemma~3.1.4}.
\end{proof}

\sourceheading[lem:3.1.7]{Lemma 3.1.7}
Let \(N\) be an irreducible \(\ell\)-adic representation of \(I\).
\begin{enumerate}[label=(\arabic*)]
\item Given any two unipotent representations \(U_1\) and \(U_2\) of \(I\),
the natural map ``\(\otimes N\)'' defines isomorphisms
\[
\begin{aligned}
  \operatorname{Hom}_I(U_1,U_2)
  &\cong
  \operatorname{Hom}_I(N\otimes U_1,N\otimes U_2),\\
  \operatorname{Ext}^1_I(U_1,U_2)
  &\cong
  \operatorname{Ext}^1_I(N\otimes U_1,N\otimes U_2).
\end{aligned}
\]
\item Given any \(\ell\)-adic representation \(M\) of \(I\) which is a
successive extension of \(N\) by itself, there exists a unipotent
representation \(U\) of \(I\) and an isomorphism
\[
  M\cong N\otimes U
\]
of \(I\)-representations. We can recover \(U\) intrinsically as the
``\(1\)-isotypical'' component of \(N^{\vee}\otimes M\).
\item Let \(M\) be an indecomposable \(\ell\)-adic representation \(M\) of
\(I\) which is a successive extension of \(N\) by itself. If \(M\) has length
\(m\) as \(\overline{\mathbb{Q}}_\ell[I]\)-module, then
\[
  M\cong N\otimes\operatorname{Unip}(m),
\]
where \(\operatorname{Unip}(m)\) is the standard Jordan block of size \(m\).
\end{enumerate}

\begin{proof}
\((1)\) Given \(U_1\) and \(U_2\) unipotent,
\[
  \operatorname{Ext}^i_I(N\otimes U_1,N\otimes U_2)
  =
  \mathrm{H}^i(I,N^{\vee}\otimes N\otimes(U_1^{\vee}\otimes U_2)).
\]
Because \(N\) is irreducible, \(N^{\vee}\otimes N\) is semisimple, and
\[
  N^{\vee}\otimes N\cong \mathbf{1}\oplus R,
\]
with \(R\cong\bigoplus\) of nontrivial irreducibles. Because
\(U_1^{\vee}\otimes U_2\) is unipotent,
\(R\otimes(U_1^{\vee}\otimes U_2)\) is a successive extension of nontrivial
irreducibles, and hence
\[
  \mathrm{H}^i(I,R\otimes(U_1^{\vee}\otimes U_2))=0
\]
for all \(i\), by \hyperref[lem:3.1.5]{Lemma~3.1.5}. Therefore the inclusion
of \(\mathbf{1}\) into \(N^{\vee}\otimes N\) induces isomorphisms
\[
  \mathrm{H}^i(I,U_1^{\vee}\otimes U_2)
  \cong
  \mathrm{H}^i(I,N^{\vee}\otimes N\otimes(U_1^{\vee}\otimes U_2))
\]
for all \(i\). If we view these groups as \(\operatorname{Ext}\) groups, then
this isomorphism is that induced by ``\(\otimes N\)''
\[
  \operatorname{Ext}^i_I(U_1,U_2)
  \longrightarrow
  \operatorname{Ext}^i_I(N\otimes U_1,N\otimes U_2)
\]
for all \(i\).

\((2)\) Once we have \((1)\), the first assertion of \((2)\) is proven by
induction on the length of \(M\). Once \(M\cong N\otimes U\), we have, just as
above,
\[
  N^{\vee}\otimes M
  =
  N^{\vee}\otimes N\otimes U
  =
  (\mathbf{1}\oplus R)\otimes U
  =
  U\oplus R\otimes U,
\]
with \(R\otimes U\) a successive extension of nontrivial irreducibles.

\((3)\) If \(N\otimes U\) is indecomposable of length \(m\), then \(U\) is
indecomposable of length \(m\), hence is \(\operatorname{Unip}(m)\).
\end{proof}

\sourceheading[prop:3.1.8]{Proposition 3.1.8}
Let \(M\) be an \(\ell\)-adic representation of \(I\). Then
\[
  \dim\left[
  \frac{\operatorname{End}_I(M)}
       {M^I\otimes(M^{\vee})^I}
  \right]
  =
  \dim\operatorname{End}_I(M/M^I).
\]

This will be proven in a series of Lemmas.

\sourceheading[lem:3.1.9]{Lemma 3.1.9}
If \hyperref[prop:3.1.8]{Proposition~3.1.8} holds for all unipotent
\(\ell\)-adic representations of \(I\), then it holds for all
\(\ell\)-adic representations of \(I\).

\begin{proof}
Let \(M\) be an \(\ell\)-adic representation of \(I\). In terms of the
``isotypical'' decomposition of \(M\) as
\(\bigoplus_{\alpha\in\mathcal{A}}M_\alpha\), we have
\[
  \operatorname{End}_I(M)=
  \bigoplus_{\alpha\in\mathcal{A}}\operatorname{End}_I(M_\alpha).
\]
Moreover, if we denote by \(\mathbf{1}\) the trivial representation of \(I\),
and by \(M_{\mathbf{1}}\) the corresponding ``isotypical'' summand (thus
\(M_{\mathbf{1}}\) is the maximal unipotent subrepresentation of \(M\)), then
clearly
\[
  M^I=(M_{\mathbf{1}})^I,\qquad
  (M^{\vee})^I=((M^{\vee})_{\mathbf{1}})^I
  =((M_{\mathbf{1}})^{\vee})^I.
\]
Thus the ``isotypical'' decomposition of \(M/M^I\) is
\[
  M/M^I
  =
  \bigoplus_{\alpha\ne\mathbf{1}\text{ in }\mathcal{A}}M_\alpha
  \oplus M_{\mathbf{1}}/(M_{\mathbf{1}})^I,
\]
and hence
\[
  \operatorname{End}_I(M/M^I)
  =
  \bigoplus_{\alpha\ne\mathbf{1}\text{ in }\mathcal{A}}\operatorname{End}_I(M_\alpha)
  \oplus
  \operatorname{End}_I(M_{\mathbf{1}}/(M_{\mathbf{1}})^I).
\]
On the other hand,
\[
  M^I\otimes(M^{\vee})^I
  =
  (M_{\mathbf{1}})^I\otimes((M^{\vee})_{\mathbf{1}})^I
  =
  (M_{\mathbf{1}})^I\otimes((M_{\mathbf{1}})^{\vee})^I,
\]
and hence
\[
\begin{aligned}
  \frac{\operatorname{End}_I(M)}{M^I\otimes(M^{\vee})^I}
  &=
  \bigoplus_{\alpha\ne\mathbf{1}\text{ in }\mathcal{A}}
  \operatorname{End}_I(M_\alpha)\\
  &\quad\oplus
  \frac{\operatorname{End}_I(M_{\mathbf{1}})}
       {(M_{\mathbf{1}})^I\otimes((M_{\mathbf{1}})^{\vee})^I}.
\end{aligned}
\]
Thus we see that \hyperref[prop:3.1.8]{Proposition~3.1.8} holds for \(M\) if
and only if it holds for \(M_{\mathbf{1}}\).
\end{proof}

\sourcepara{3.1.10}
We now study the case in which \(M\) is a unipotent \(\ell\)-adic
representation \(U\) of \(I\), so a direct sum of standard Jordan blocks
\(\bigoplus\operatorname{Unip}(n_i)\) of varying sizes \(n_i\geq 1\). Given a
unipotent representation \(U\) of \(I\), we define a sequence of nonnegative
integers \(e_i=e_i(U)\),
\[
  e_1\geq e_2\geq e_3\geq e_4\geq\cdots,\qquad e_r=0\text{ for }r\gg 0,
\]
by
\[
  e_i:=\text{the number of Jordan blocks of dimension }\geq i.
\]
It is obvious from this definition that given two unipotent representations
\(U_1\) and \(U_2\) of \(I\), we have
\[
  e_i(U_1\oplus U_2)=e_i(U_1)+e_i(U_2).
\]
For a single unipotent block \(\operatorname{Unip}(n)\), we have
\[
\begin{aligned}
  e_i&=1 &&\text{for }i\leq n,\\
  e_i&=0 &&\text{for }i>n.
\end{aligned}
\]
So for \(M=\bigoplus\operatorname{Unip}(n_i)\), the \(e_i(M)\) are the
partition of \(\dim(M)\), written in decreasing order, which is dual to the
partition of \(\dim(M)\) given by the block sizes \(n_i\). We call the sequence
\((e_1,e_2,e_3,\ldots)\) attached to a unipotent \(M\) its ``dual partition''.

\sourcepara{3.1.11}
Clearly if \(M\) is unipotent, with dual partition
\((e_1,e_2,e_3,e_4,\ldots)\), then \(M/M^I\) has dual partition
\((e_2,e_3,e_4,\ldots)\), formed from that of \(M\) by discarding the leading
term. (Indeed, if \(M=\operatorname{Unip}(n)\), then \(M/M^I\) is \(0\) if
\(n=1\), and is \(\operatorname{Unip}(n-1)\) if \(n>1\).)

\sourcepara{3.1.12}
Notice that if \(M\) is a single Jordan block \(\operatorname{Unip}(n)\), then
\(M^{\vee}\) is also a single Jordan block \(\operatorname{Unip}(n)\). So if
\(M\cong\bigoplus\operatorname{Unip}(n_i)\), then also
\(M^{\vee}\cong\bigoplus\operatorname{Unip}(n_i)\). From this we see that
\[
\begin{aligned}
  \dim(M^I)&=\text{the number of Jordan blocks in }M:=e_1(M),\\
  \dim((M^{\vee})^I)&=e_1(M^{\vee})=e_1(M).
\end{aligned}
\]

\sourceheading[rem:3.1.13]{Remark 3.1.13}
One could define the dual partition inductively by the properties
\[
  e_1(M)=\dim(M^I),\qquad e_{i+1}(M)=e_i(M/M^I),\quad\text{for }i\geq 1.
\]

\sourceheading[lem:3.1.14]{Lemma 3.1.14}
If \(M\) and \(N\) are unipotent \(\ell\)-adic representations of \(I\), with
dual partitions \((e_1,e_2,e_3,e_4,\ldots)\) and
\((f_1,f_2,f_3,f_4,\ldots)\), then
\[
  \dim\operatorname{Hom}_I(M,N)=\sum_i e_i f_i.
\]

\begin{proof}
Both sides are additive over direct sums, so it suffices to check in the case
when \(M\) is \(\operatorname{Unip}(m)\) and \(N\) is
\(\operatorname{Unip}(n)\). In this case the assertion is that
\[
  \dim\operatorname{Hom}_I(\operatorname{Unip}(m),\operatorname{Unip}(n))
  =
  \min(m,n).
\]
In terms of the action of
\[
  T:=(\text{a topological generator of }I/P)-1,
\]
this is the assertion that over any field \(E\), we have
\[
  \dim_E\operatorname{Hom}_{E[T]}(E[T]/(T^m),E[T]/(T^n))=\min(m,n).
\]
But by the map \(\varphi\mapsto\varphi(1)\),
\[
\begin{aligned}
  \operatorname{Hom}_{E[T]}(E[T]/(T^m),E[T]/(T^n))
  &\cong \text{kernel of }T^m\text{ in }E[T]/(T^n)\\
  &=E[T]/(T^n), &&\text{if }m\geq n,\\
  &=T^{n-m}E[T]/(T^n)\cong E[T]/(T^m), &&\text{if }n\geq m.
\end{aligned}
\]
\end{proof}

\sourceheading[lem:3.1.15]{Lemma 3.1.15}
If \(M\) is a unipotent \(\ell\)-adic representation of \(I\), with dual
partition \((e_1,e_2,e_3,e_4,\ldots)\), then
\[
  \dim\operatorname{End}_I(M)=\sum_i(e_i)^2.
\]

\begin{proof}
Take \(N=M\) in \hyperref[lem:3.1.14]{Lemma~3.1.14}.
\end{proof}

\sourceheading[lem:3.1.16]{Lemma 3.1.16}
If \(M\) is a unipotent \(\ell\)-adic representation of \(I\),
\[
  \dim\operatorname{End}_I(M/M^I)
  =
  \dim\left[
  \frac{\operatorname{End}_I(M)}
       {M^I\otimes(M^{\vee})^I}
  \right].
\]

\begin{proof}
If \(M\) has dual partition \((e_1,e_2,e_3,e_4,\ldots)\), then \(M/M^I\) has
dual partition \((e_2,e_3,e_4,\ldots)\), and
\(\dim(M^I\otimes(M^{\vee})^I)=(e_1)^2\), so both sides are equal to
\(\sum_{i\geq 2}(e_i)^2\).
\end{proof}

In view of \hyperref[lem:3.1.9]{Lemma~3.1.9}, this completes the proof of
\hyperref[prop:3.1.8]{Proposition~3.1.8}.

\section{Interlude: the operation
\texorpdfstring{\(\otimes_{\mathrm{mid}}\)}{tensor-mid}}\label{sec:3.2}

\sourcepara{3.2.1}
Let \(k\) be a field, \(X\) a separated \(k\)-scheme of finite type which is
irreducible and smooth, everywhere of relative dimension \(n\). Let \(\ell\)
be a prime number \(\ne\operatorname{char}(k)\). We look at the full
subcategory \(\operatorname{ME}(X)\) of \((\ell\)-adic) \(\operatorname{Perv}(X)\)
consisting of those objects which are middle extensions from every nonempty
affine open set, i.e., those perverse objects \(K\) on \(X\) such that for
every nonempty affine open set \(j:U\to X\), we have
\[
  K\cong j_{!*}j^*K.
\]
Given \(K\) in \(\operatorname{ME}\), we say that \(K\) has generic rank
\(r\) if for some (or equivalently, for any) nonempty affine open set
\(j:U\to X\) on which it is lisse, \(j^*K\) is \(\mathcal{F}[n]\), with
\(\mathcal{F}\) a lisse sheaf on \(U\) of rank \(r\).

\sourcepara{3.2.2}
We define an operation
\[
  \otimes_{\mathrm{mid}}:\operatorname{ME}(X)\times\operatorname{ME}(X)
  \longrightarrow \operatorname{ME}(X)
\]
as follows. Given any two objects \(K\) and \(L\) in \(\operatorname{ME}\), pick
a common nonempty affine open set \(j:U\to X\) on which both \(K\) and \(L\)
are lisse. Thus \(j^*K=\mathcal{F}[n]\) and \(j^*L=\mathcal{G}[n]\) for lisse
sheaves \(\mathcal{F}\) and \(\mathcal{G}\) on \(U\). We define
\[
  K\otimes_{\mathrm{mid}}L
  :=
  j_{!*}\bigl((\mathcal{F}\otimes\mathcal{G})[n]\bigr)
  =
  j_{!*}\bigl((j^*K\otimes j^*L)[-n]\bigr).
\]
This is easily seen to be independent of the auxiliary choice of \(U\), using
\hyperref[par:2.3.3.1]{(2.3.3.1)}.

\sourceheading[lem:3.2.3]{Lemma 3.2.3}
Let \(k\) be a field, \(X\) a separated \(k\)-scheme of finite type which is
irreducible and smooth, everywhere of relative dimension \(n\). Let \(\ell\)
be a prime number \(\ne\operatorname{char}(k)\).
\begin{enumerate}[label=(\arabic*)]
\item An object \(K\) in \(\operatorname{ME}\) is perverse irreducible if and
only if it is irreducible in \(\operatorname{ME}\).
\item Let \(L\) in \(\operatorname{ME}\) have generic rank one. Then
\begin{enumerate}[label=(2\alph*)]
\item The operation \(K\mapsto K\otimes_{\mathrm{mid}}L\) is an
autoequivalence of \(\operatorname{ME}\) with
itself, whose inverse is \(K\mapsto K\otimes_{\mathrm{mid}}\mathbb{D}L\).
\item The operation \(K\mapsto K\otimes_{\mathrm{mid}}L\) induces an
automorphism of the family of all irreducible perverse sheaves in
\(\operatorname{ME}\), whose inverse is
\(K\mapsto K\otimes_{\mathrm{mid}}\mathbb{D}L\).
\end{enumerate}
\end{enumerate}

\begin{proof}
\((1)\) Since \(\operatorname{ME}\) is a full subcategory of
\(\operatorname{Perv}\), any perverse irreducible is an irreducible object of
\(\operatorname{ME}\). Conversely, suppose that \(K\) in \(\operatorname{ME}\)
is not perverse irreducible. Then there exists a perverse irreducible \(L\),
and a nonzero element in \(\operatorname{Hom}_X(L,K)\). If \(L\) is in
\(\operatorname{ME}\), \(K\) is not irreducible in \(\operatorname{ME}\). If
not, then \(L\) is supported on a proper closed subvariety \(Y\) of \(X\). So
there exists an affine open \(j:U\to X\) with \(j^*L=0\). But in
\(\operatorname{Perv}(X)\), we have \(K=j_{!*}j^*K\subset \mathrm{R}j_*j^*K\),
so
\[
  \operatorname{Hom}_X(L,K)
  \subset
  \operatorname{Hom}_X(L,\mathrm{R}j_*j^*K)
  =
  \operatorname{Hom}_U(j^*L,j^*K)
  =
  0,
\]
contradiction. \((2a)\) is obvious by \hyperref[par:2.3.3.1]{(2.3.3.1)}, and
by \((1)\) it implies \((2b)\).
\end{proof}

\sourceheading[lem:3.2.4]{Lemma 3.2.4}
Let \(\ell\) be a prime number. On \(\mathbb{A}^1\) over an algebraically
closed field of characteristic \(\ne\ell\), let \(L\) in
\(\mathrm{D}^b_c(\mathbb{A}^1,\overline{\mathbb{Q}}_\ell)\) be a perverse
sheaf on \(\mathbb{A}^1\) which lies in \(\operatorname{ME}\). Suppose that
\(L\) has generic rank one. Then
\begin{enumerate}[label=(\arabic*)]
\item If \(K\) in \(\operatorname{ME}\) is perverse irreducible, so is
\(K\otimes_{\mathrm{mid}}L\).
\item For any \(K\) in \(\operatorname{ME}\),
\[
  \operatorname{rig}(K)=\operatorname{rig}(K\otimes_{\mathrm{mid}}L).
\]
\end{enumerate}

\begin{proof}
Assertion \((1)\) is part \((2b)\) of
\hyperref[lem:3.2.3]{Lemma~3.2.3}. Assertion \((2)\) is obvious from the
definitions: if \(j:U\to\mathbb{A}^1\) is any nonempty open set on which both
\(K\) and \(L\) are lisse, say
\(j^*K=\mathcal{F}[1]\) and \(j^*L=\mathcal{L}[1]\), with \(\mathcal{F}\)
lisse of some rank \(r\), and \(\mathcal{L}\) lisse of rank one, then on \(U\)
the sheaves \(\operatorname{End}(\mathcal{F})\) and
\(\operatorname{End}(\mathcal{F}\otimes\mathcal{L})\) coincide.
\end{proof}

\section{Applications to middle additive convolution}\label{sec:3.3}

\sourceheading[lem:3.3.1]{Lemma 3.3.1}
Let \(p\) and \(\ell\) be prime numbers, \(\ell\ne p\). On \(\mathbb{A}^1\)
over an algebraically closed field of characteristic \(p\), let \(K\) and
\(L\) in \(\mathrm{D}^b_c(\mathbb{A}^1,\overline{\mathbb{Q}}_\ell)\) be
perverse sheaves. Then
\begin{enumerate}[label=(\arabic*)]
\item \(K\) is in \(\mathcal{P}\) if and only if \(\operatorname{FT}K\) is in
\(\operatorname{ME}\).
\item If \(K\) and \(L\) are both in \(\mathcal{P}\), then
\[
  \operatorname{FT}(K*_{\mathrm{mid}+}L)
  =
  \operatorname{FT}K\otimes_{\mathrm{mid}}\operatorname{FT}L.
\]
\end{enumerate}

\begin{proof}
This was proven in \hyperref[lem:2.10.2]{Lemma~2.10.2} and
\hyperref[thm:2.10.8]{Theorem~2.10.8} above.
\end{proof}

\sourceheading[lem:3.3.2]{Lemma 3.3.2}
Let \(p\) and \(\ell\) be prime numbers, \(\ell\ne p\). On \(\mathbb{G}_m\)
over an algebraically closed field of characteristic \(p>0\), let \(\chi\) be
any continuous nontrivial \(\overline{\mathbb{Q}}_\ell^\times\)-valued
character of \(\pi_1(\mathbb{G}_m)^{\mathrm{tame}}\), \(\mathcal{L}_\chi\) the
corresponding Kummer sheaf on \(\mathbb{G}_m\), and \(j_*\mathcal{L}_\chi\)
its extension by direct image to \(\mathbb{A}^1\). Then
\[
  \operatorname{FT}(j_*\mathcal{L}_{\chi}[1])
  =
  j_*\mathcal{L}_{\overline{\chi}}[1]
\]
(geometrically).

\begin{proof}
Direct calculation.
\end{proof}

\sourceheading[thm:3.3.3]{Theorem 3.3.3}
Let \(p\) and \(\ell\) be prime numbers, \(\ell\ne p\). Over an algebraically
closed field of characteristic \(p\), let \(\chi\) be any continuous nontrivial
\(\overline{\mathbb{Q}}_\ell^\times\)-valued character of
\(\pi_1(\mathbb{G}_m)^{\mathrm{tame}}\). Let \(K\) in
\(\mathrm{D}^b_c(\mathbb{A}^1,\overline{\mathbb{Q}}_\ell)\) be perverse
irreducible on \(\mathbb{A}^1\).
\begin{enumerate}[label=(\arabic*)]
\item If \(K\) is not in \(\mathcal{P}\), then
\begin{enumerate}[label=(1\alph*)]
\item If \(K\) is \(\overline{\mathbb{Q}}_\ell[1]\), then
\[
  K*_{\mathrm{mid}+}j_*\mathcal{L}_{\chi}[1]=0.
\]
\item If \(K\) is \(\mathcal{L}_{\psi_\alpha}[1]\) with \(\alpha\ne0\), then
\[
  K*_{\mathrm{mid}+}j_*\mathcal{L}_{\chi}[1]
  =
  \mathcal{L}_{\psi_\alpha}[1].
\]
\end{enumerate}
\item If \(K\) is in \(\mathcal{P}\), then
\begin{enumerate}[label=(2\alph*)]
\item If \(K\) is punctual, say \(\delta_\alpha\),
\[
  K*_{\mathrm{mid}+}j_*\mathcal{L}_{\chi}[1]
  =
  j_*\mathcal{L}_{\chi}(x-\alpha)[1].
\]
\item If \(K\) is \(j_*\mathcal{L}_{\overline{\chi}}(x-\alpha)[1]\) for some
\(\alpha\), then
\[
  K*_{\mathrm{mid}+}j_*\mathcal{L}_{\chi}[1]=\delta_\alpha.
\]
\item If \(K\) is \(j_*\mathcal{L}_{\rho}(x-\alpha)[1]\) for some \(\alpha\)
and some \(\rho\ne\overline{\chi}\), then
\[
  K*_{\mathrm{mid}+}j_*\mathcal{L}_{\chi}[1]
  =
  j_*\mathcal{L}_{\rho\chi}(x-\alpha)[1].
\]
\item If \(K\) is a perverse irreducible in \(\mathcal{P}\) which is not of
type \((2a)\), \((2b)\), or \((2c)\), then
\[
  K*_{\mathrm{mid}+}j_*\mathcal{L}_{\chi}[1]
\]
is a perverse irreducible in \(\mathcal{P}\) which is not of type \((2a)\),
\((2b)\), or \((2c)\).
\end{enumerate}
\item If \(K\) is in \(\mathcal{P}\) and of type \((2d)\), then
\[
  \operatorname{rig}(K)
  =
  \operatorname{rig}(K*_{\mathrm{mid}+}j_*\mathcal{L}_{\chi}[1]).
\]
\end{enumerate}

\begin{proof}
Assertion \((1)\) is proven by direct calculation. Assertion \((2a)\) is
proven by direct calculation. We know that middle convolution
\(*_{\mathrm{mid}+}j_*\mathcal{L}_{\chi}[1]\) with
\(j_*\mathcal{L}_{\chi}[1]\) is an automorphism of the irreducibles in
\(\mathcal{P}\), with inverse
\(*_{\mathrm{mid}+}j_*\mathcal{L}_{\overline{\chi}}[1]\). Using this,
\((2b)\) follows from \((2a)\). We also know that
\[
\begin{aligned}
  j_*\mathcal{L}_{\chi}[1]*_{\mathrm{mid}+}j_*\mathcal{L}_{\rho}[1]
  &=
  j_*\mathcal{L}_{\chi\rho}[1] &&\text{if }\chi\rho\ne\mathbf{1},\\
  &=\delta_0 &&\text{if }\chi\rho=\mathbf{1}.
\end{aligned}
\]
Using the associativity of middle convolution, and \((2a)\), we get \((2c)\).

To prove \((2d)\), observe first that for each \(\alpha\), the family of all
those irreducibles in \(\mathcal{P}\) which are either \(\delta_\alpha\) or
\(j_*\mathcal{L}_\rho(x-\alpha)[1]\) for some \(\rho\ne\mathbf{1}\) is simply
the orbit of \(\delta_\alpha\) under the group of operators given by middle
convolution with \(\delta_0\) and with any \(j_*\mathcal{L}_{\chi}[1]\), for
any \(\chi\ne\mathbf{1}\). Since the family of all irreducibles in
\(\mathcal{P}\) is also stable by this group of operators, the complement of
a union of orbits is also stable.

To prove \((3)\), we note that by \((2d)\), both sides are middle extensions.
Since both sides are also in \(\mathcal{P}\), their Fourier Transforms are
middle extensions. So by \hyperref[thm:3.0.2]{Theorem~3.0.2}, we have
\[
  \operatorname{rig}(K)=\operatorname{rig}(\operatorname{FT}K),
\]
and
\[
\begin{aligned}
  \operatorname{rig}(K*_{\mathrm{mid}+}j_*\mathcal{L}_{\chi}[1])
  &=
  \operatorname{rig}\bigl(\operatorname{FT}
  (K*_{\mathrm{mid}+}j_*\mathcal{L}_{\chi}[1])\bigr)\\
  &=
  \operatorname{rig}\bigl(\operatorname{FT}K\otimes_{\mathrm{mid}}
  \operatorname{FT}(j_*\mathcal{L}_{\chi}[1])\bigr),
  \quad\text{by \hyperref[lem:3.3.1]{Lemma~3.3.1}(2)}\\
  &=
  \operatorname{rig}\bigl(\operatorname{FT}K\otimes_{\mathrm{mid}}
  j_*\mathcal{L}_{\overline{\chi}}[1]\bigr),
  \quad\text{by \hyperref[lem:3.3.2]{Lemma~3.3.2}}\\
  &=
  \operatorname{rig}(\operatorname{FT}K),
  \quad\text{by \hyperref[lem:3.2.4]{Lemma~3.2.4}(2)}.
\end{aligned}
\]
\end{proof}

\sourcepara{3.3.4}
We now turn to a discussion of the monodromy of middle convolution with a
\(j_*\mathcal{L}_{\chi}[1]\). For the convenience of the reader, we recall
some of the basic facts about Laumon's local Fourier Transforms.

\sourceheading[thm:3.3.5]{Theorem 3.3.5}
Hypotheses and notations as in \hyperref[thm:3.3.3]{Theorem~3.3.3} above,
suppose that \(K\) is a perverse irreducible in \(\mathcal{P}\) of type
\((2d)\). Pick a nonempty open set \(j:U\to\mathbb{A}^1\) on which both \(K\)
and \(K*_{\mathrm{mid}+}j_*\mathcal{L}_{\chi}[1]\) are lisse, say
\[
  K=j_{!*}\mathcal{F}[1],\qquad
  K*_{\mathrm{mid}+}j_*\mathcal{L}_{\chi}[1]
  =
  j_{!*}\mathcal{G}[1],
\]
with \(\mathcal{F}\) and \(\mathcal{G}\) lisse sheaves on \(U\). Then the local
monodromies of \(\mathcal{F}\) and of \(\mathcal{G}\) are related as follows:
\begin{enumerate}[label=(\arabic*)]
\item at \(s\) in \(\mathbb{A}^1-U\):
\[
  \operatorname{FT}_{\psi}^{\mathrm{loc}}(s,\infty)
  \bigl(\mathcal{F}(s)/\mathcal{F}(s)^{I(s)}\bigr)
  \otimes\mathcal{L}_{\overline{\chi}}
  =
  \operatorname{FT}_{\psi}^{\mathrm{loc}}(s,\infty)
  \bigl(\mathcal{G}(s)/\mathcal{G}(s)^{I(s)}\bigr).
\]
\item[(1a)] In particular, \(\mathcal{F}\) is lisse at \(s\) if and only if
\(\mathcal{G}\) is lisse at \(s\).
\item at \(\infty\): decompose
\[
  \mathcal{F}(\infty)
  =
  \mathcal{F}(\infty)(\text{slopes }>1)
  \oplus
  \left(\bigoplus_{s\in\mathbb{A}^1}
  \mathcal{L}_{\psi_s}\otimes\mathcal{F}(\infty,s)\right)
\]
with \(\mathcal{F}(\infty,s)\) an \(I(\infty)\)-representation with all slopes
\(<1\), and similarly for \(\mathcal{G}\). Then
\item[(2a)] We have the formula
\[
  \operatorname{FT}_{\psi}^{\mathrm{loc}}(\infty,\infty)
  \bigl(\mathcal{F}(\infty)(\text{slopes }>1)\bigr)
  \otimes\mathcal{L}_{\overline{\chi}}
  =
  \operatorname{FT}_{\psi}^{\mathrm{loc}}(\infty,\infty)
  \bigl(\mathcal{G}(\infty)(\text{slopes }>1)\bigr).
\]
\item[(2b)] For each \(s\ne0\) in \(\mathbb{A}^1\),
\[
  \mathcal{F}(\infty,s)\cong\mathcal{G}(\infty,s).
\]
\item[(2c)] There exists an \(I(0)\)-representation \(M(0)\) with
\[
\begin{aligned}
  \mathcal{F}(\infty,0)
  &=
  \operatorname{FT}_{\overline{\psi}}^{\mathrm{loc}}(0,\infty)
  \bigl(M(0)/M(0)^{I(0)}\bigr),\\
  \mathcal{G}(\infty,0)
  &=
  \operatorname{FT}_{\overline{\psi}}^{\mathrm{loc}}(0,\infty)
  \bigl(M(0)\otimes\mathcal{L}_{\overline{\chi}}
  /(M(0)\otimes\mathcal{L}_{\overline{\chi}})^{I(0)}\bigr).
\end{aligned}
\]
\item[(2d)] In particular, cf. \cite[7.4.4]{Ka-ESDE}, \(\mathcal{F}\) is
tame at \(\infty\) if and only if \(\mathcal{G}\) is tame at \(\infty\).
\end{enumerate}

\begin{proof}
This is just the spelling out of identities
\[
\begin{aligned}
  \operatorname{FT}_{\psi}(K*_{\mathrm{mid}+}L)
  &=
  \operatorname{FT}_{\psi}K\otimes_{\mathrm{mid}}\operatorname{FT}_{\psi}L,\\
  \operatorname{FT}_{\overline{\psi}}
  (\operatorname{FT}_{\psi}K\otimes_{\mathrm{mid}}\operatorname{FT}_{\psi}L)
  &=
  K*_{\mathrm{mid}+}L,\\
  \operatorname{FT}_{\overline{\psi}}(\operatorname{FT}_{\psi}K)&=K,
\end{aligned}
\]
with \(L=j_*\mathcal{L}_{\chi}[1]\), together with the effect of
\(\operatorname{FT}\) upon local monodromies, as expressed via Laumon's
\(\operatorname{FTloc}(s,\infty)\) and \(\operatorname{FTloc}(\infty,\infty)\)
functors.

Concretely, let us write
\(\operatorname{FT}_{\psi}K=j_{!*}(\mathcal{A}[1])\), for some sufficiently
small nonempty open set \(j:V\to\mathbb{A}^1\) and some lisse sheaf
\(\mathcal{A}\) on \(V\). Then
\(\operatorname{FT}_{\psi}(K*_{\mathrm{mid}+}L)=j_{!*}(\mathcal{B}[1])\), with
\(\mathcal{B}=\mathcal{A}\otimes j^*\mathcal{L}_{\overline{\chi}}\). Their
\(I(\infty)\)-representations are related by
\(\mathcal{B}(\infty)=\mathcal{A}(\infty)\otimes\mathcal{L}_{\overline{\chi}}\),
so decomposing them we find
\[
\begin{aligned}
  \mathcal{B}(\infty)(\text{slope }>1)
  &=
  \mathcal{A}(\infty)(\text{slope }>1)\otimes
  \mathcal{L}_{\overline{\chi}},\\
  \mathcal{B}(\infty,s)
  &=
  \mathcal{A}(\infty,s)\otimes\mathcal{L}_{\overline{\chi}},
  \quad\text{for every }s\text{ in }\mathbb{A}^1.
\end{aligned}
\]
By Laumon (cf. \cite[7.4.2]{Ka-ESDE}) we know that
\[
\begin{aligned}
  \operatorname{FT}_{\psi}^{\mathrm{loc}}(s,\infty)
  \bigl(\mathcal{F}(s)/\mathcal{F}(s)^{I(s)}\bigr)
  &=
  \mathcal{A}(\infty,s),\\
  \operatorname{FT}_{\psi}^{\mathrm{loc}}(s,\infty)
  \bigl(\mathcal{G}(s)/\mathcal{G}(s)^{I(s)}\bigr)
  &=
  \mathcal{B}(\infty,s)
  =
  \mathcal{A}(\infty,s)\otimes\mathcal{L}_{\overline{\chi}},
\end{aligned}
\]
whence
\[
  \operatorname{FT}_{\psi}^{\mathrm{loc}}(s,\infty)
  \bigl(\mathcal{F}(s)/\mathcal{F}(s)^{I(s)}\bigr)
  \otimes\mathcal{L}_{\overline{\chi}}
  =
  \operatorname{FT}_{\psi}^{\mathrm{loc}}(s,\infty)
  \bigl(\mathcal{G}(s)/\mathcal{G}(s)^{I(s)}\bigr).
\]
By Laumon, we also know
\[
\begin{aligned}
  \operatorname{FT}_{\psi}^{\mathrm{loc}}(\infty,\infty)
  \bigl(\mathcal{F}(\infty)(\text{slope }>1)\bigr)
  &=
  \mathcal{A}(\infty)(\text{slope }>1),\\
  \operatorname{FT}_{\psi}^{\mathrm{loc}}(\infty,\infty)
  \bigl(\mathcal{G}(\infty)(\text{slope }>1)\bigr)
  &=
  \mathcal{B}(\infty)(\text{slope }>1)\\
  &=
  \mathcal{A}(\infty)(\text{slope }>1)\otimes
  \mathcal{L}_{\overline{\chi}},
\end{aligned}
\]
whence the asserted formula for the slope \(>1\) part.

To prove the remaining assertions, we use the inversion formulas
\[
\begin{aligned}
  \operatorname{FT}_{\overline{\psi}}
  (\operatorname{FT}_{\psi}K\otimes_{\mathrm{mid}}\operatorname{FT}_{\psi}L)
  &=
  K*_{\mathrm{mid}+}L,\\
  \operatorname{FT}_{\overline{\psi}}(\operatorname{FT}_{\psi}K)&=K.
\end{aligned}
\]
Since our \(\operatorname{FT}L\) is lisse of rank one on \(\mathbb{G}_m\),
\(\operatorname{FT}K\) and \(\operatorname{FT}K\otimes_{\mathrm{mid}}
\operatorname{FT}L\) give isomorphic representations of \(I(s)\) for each
\(s\ne0\) in \(\mathbb{A}^1\):
\[
  \mathcal{A}(s)\cong\mathcal{B}(s)\quad
  \text{for }s\ne0\text{ in }\mathbb{A}^1.
\]
By stationary phase,
\[
\begin{aligned}
  \mathcal{F}(\infty,s)
  &=
  \operatorname{FT}_{\overline{\psi}}^{\mathrm{loc}}(s,\infty)
  \bigl(\mathcal{A}(s)/\mathcal{A}(s)^{I(s)}\bigr),\\
  \mathcal{G}(\infty,s)
  &=
  \operatorname{FT}_{\overline{\psi}}^{\mathrm{loc}}(s,\infty)
  \bigl(\mathcal{B}(s)/\mathcal{B}(s)^{I(s)}\bigr),
  \quad\text{for every }s\text{ in }\mathbb{A}^1.
\end{aligned}
\]
But \(\mathcal{A}(s)\cong\mathcal{B}(s)\) for \(s\ne0\), so
\(\mathcal{F}(\infty,s)\cong\mathcal{G}(\infty,s)\) for \(s\ne0\). For \(s=0\),
\[
  \mathcal{A}(0)\otimes\mathcal{L}_{\overline{\chi}}\cong\mathcal{B}(0),
\]
so with \(M(0):=\mathcal{A}(0)\) we get
\[
\begin{aligned}
  \mathcal{F}(\infty,0)
  &=
  \operatorname{FT}_{\overline{\psi}}^{\mathrm{loc}}(0,\infty)
  \bigl(M(0)/M(0)^{I(0)}\bigr),\\
  \mathcal{G}(\infty,0)
  &=
  \operatorname{FT}_{\overline{\psi}}^{\mathrm{loc}}(0,\infty)
  \bigl(M(0)\otimes\mathcal{L}_{\overline{\chi}}
  /(M(0)\otimes\mathcal{L}_{\overline{\chi}})^{I(0)}\bigr).
\end{aligned}
\]
\end{proof}

\sourceheading[cor:3.3.6]{Corollary 3.3.6}
Hypotheses and notations as in \hyperref[thm:3.3.3]{Theorem~3.3.3} above,
suppose in addition that \(K\) is tamely ramified everywhere. Then so is
\(K*_{\mathrm{mid}+}j_*\mathcal{L}_{\chi}[1]\), and their local monodromies are
related as follows:
\begin{enumerate}[label=(\arabic*)]
\item at \(s\) in \(\mathbb{A}^1-U\):
\[
  \bigl(\mathcal{F}(s)/\mathcal{F}(s)^{I(s)}\bigr)
  \otimes\mathcal{L}_{\chi(x-s)}
  \cong
  \mathcal{G}(s)/\mathcal{G}(s)^{I(s)}.
\]
\item at \(\infty\): there exists a tame \(I(\infty)\)-representation
\(M(\infty)\) with
\[
\begin{aligned}
  \mathcal{F}(\infty)
  &=
  M(\infty)/M(\infty)^{I(\infty)},\\
  \mathcal{G}(\infty)
  &=
  M(\infty)\otimes\mathcal{L}_{\chi}/
  (M(\infty)\otimes\mathcal{L}_{\chi})^{I(\infty)}.
\end{aligned}
\]
\item \(M(\infty)\) is the unique \(I(\infty)\)-representation with these
properties.
\item We have the formulas
\[
\begin{aligned}
  \operatorname{rank}M(\infty)
  &=
  \sum_{s\in\mathbb{A}^1-U}
  \operatorname{rank}\bigl(\mathcal{F}(s)/\mathcal{F}(s)^{I(s)}\bigr),\\
  \operatorname{rank}M(\infty)
  &=
  \sum_{s\in\mathbb{A}^1-U}
  \operatorname{rank}\bigl(\mathcal{G}(s)/\mathcal{G}(s)^{I(s)}\bigr).
\end{aligned}
\]
\end{enumerate}

\begin{proof}
At \(s\) in \(\mathbb{A}^1-U\), we have
\[
  \operatorname{FT}_{\psi}^{\mathrm{loc}}(s,\infty)
  \bigl(\mathcal{F}(s)/\mathcal{F}(s)^{I(s)}\bigr)
  \otimes\mathcal{L}_{\overline{\chi}}
  =
  \operatorname{FT}_{\psi}^{\mathrm{loc}}(s,\infty)
  \bigl(\mathcal{G}(s)/\mathcal{G}(s)^{I(s)}\bigr).
\]
Since \(\mathcal{F}(s)\) is tame, and
\(\operatorname{FT}_{\psi}^{\mathrm{loc}}(s,\infty)\) and its quasi-inverse
both carry tames to tames, we see that
\(\mathcal{G}(s)/\mathcal{G}(s)^{I(s)}\) is tame, whence \(\mathcal{G}(s)\)
itself is tame. Moreover, for \(M\) a tame \(I(s)\)-representation, we have
\cite[7.4.1.3]{Ka-ESDE}
\[
\begin{aligned}
  \operatorname{FT}_{\psi}^{\mathrm{loc}}(s,\infty)(M)
  \otimes\mathcal{L}_{\overline{\chi}}
  &=
  \operatorname{FT}_{\psi}^{\mathrm{loc}}(s,\infty)
  (M\otimes\mathcal{L}_{\chi(x-s)}),\\
  \operatorname{FT}_{\psi}^{\mathrm{loc}}(s,\infty)(M)
  &=
  [x-s\mapsto 1/(x-s)]^*M
\end{aligned}
\]
as \(I(\infty)\)-representations. So from the above isomorphism of
\(I(\infty)\)-representations we infer
\[
  \bigl(\mathcal{F}(s)/\mathcal{F}(s)^{I(s)}\bigr)
  \otimes\mathcal{L}_{\chi(x-s)}
  =
  \mathcal{G}(s)/\mathcal{G}(s)^{I(s)}.
\]
We now turn to the situation at \(\infty\). Since \(K\) is tame at \(\infty\),
we have \(\mathcal{F}(\infty)(\text{slopes }>1)=0\). Therefore
\(\operatorname{FT}_{\psi}^{\mathrm{loc}}(\infty,\infty)
(\mathcal{G}(\infty)(\text{slopes }>1))=0\), and hence
\(\mathcal{G}(\infty)(\text{slopes }>1)=0\). For \(s\ne0\),
\(\mathcal{F}(\infty,s)=0\), hence \(\mathcal{G}(\infty,s)=0\). For \(s=0\),
any \(I(0)\)-representation \(M(0)\) with
\[
  \mathcal{F}(\infty,0)
  =
  \operatorname{FT}_{\overline{\psi}}^{\mathrm{loc}}(0,\infty)
  \bigl(M(0)/M(0)^{I(0)}\bigr)
\]
must be tame \cite[7.4.4.4]{Ka-ESDE}. Moreover, for tame
\(I(0)\)-representations, we have \cite[7.4.4.3]{Ka-ESDE}
\[
  \operatorname{FT}_{\overline{\psi}}^{\mathrm{loc}}(0,\infty)(M(0))
  =
  [x\mapsto 1/x]^*M(0)
\]
as \(I(\infty)\)-representations. We take our \(M(\infty)\) to be
\([x\mapsto 1/x]^*M(0)\) for \(M(0)\) appearing in
\hyperref[thm:3.3.5]{Theorem~3.3.5} (i.e., if
\(\operatorname{FT}_{\psi}K=j_{!*}\mathcal{A}[1]\), for a sufficiently small
nonempty open set \(j:V\to\mathbb{A}^1\) and a lisse sheaf \(\mathcal{A}\) on
\(V\), then \(M(0)\) is \(\mathcal{A}(0)\)).

To show that \(M(\infty)\) is uniquely determined as \(I(\infty)\)-
representation by the two conditions
\[
\begin{aligned}
  \mathcal{F}(\infty)&=M(\infty)/M(\infty)^{I(\infty)},\\
  \mathcal{G}(\infty)&=
  M(\infty)\otimes\mathcal{L}_{\chi}/
  (M(\infty)\otimes\mathcal{L}_{\chi})^{I(\infty)},
\end{aligned}
\]
write the ``isotypical decomposition'' of \(M(\infty)\), cf.
\hyperref[lem:3.1.6]{Lemma~3.1.6}, say
\[
  M(\infty)=\bigoplus_{\alpha\in\mathcal{A}}M(\infty)_\alpha.
\]
By the first condition,
\[
  M(\infty)_\alpha\cong\mathcal{F}(\infty)_\alpha
  \quad\text{for }\alpha\ne\mathbf{1}.
\]
By the second,
\[
  (M(\infty)\otimes\mathcal{L}_{\chi})_\alpha
  \cong
  \mathcal{G}(\infty)_\alpha\quad\text{for }\alpha\ne\mathbf{1},
\]
which we rewrite as
\[
  (M(\infty)\otimes\mathcal{L}_{\chi})_{\alpha\otimes\chi}
  \cong
  \mathcal{G}(\infty)_{\alpha\otimes\chi}
  \quad\text{for }\alpha\otimes\chi\ne\mathbf{1},
\]
i.e.,
\[
  M(\infty)_\alpha
  \cong
  (\mathcal{G}(\infty)\otimes\mathcal{L}_{\overline{\chi}})_\alpha
  \quad\text{for }\alpha\ne\overline{\chi}.
\]
In particular,
\[
  M(\infty)_{\mathbf{1}}
  \cong
  (\mathcal{G}(\infty)\otimes\mathcal{L}_{\overline{\chi}})_{\mathbf{1}}.
\]
Thus
\[
  M(\infty)
  =
  \bigoplus_{\alpha\ne\mathbf{1}}\mathcal{F}(\infty)_\alpha
  \oplus
  (\mathcal{G}(\infty)\otimes\mathcal{L}_{\overline{\chi}})_{\mathbf{1}}.
\]
To compute the rank of \(M(\infty)\), recall that if we write
\(\operatorname{FT}_{\psi}K=j_{!*}(\mathcal{A}[1])\), for a sufficiently small
nonempty open set \(j:V\to\mathbb{A}^1\) and a lisse sheaf \(\mathcal{A}\) on
\(V\), then \(M(\infty)\) is \([x\mapsto 1/x]^*\mathcal{A}(0)\). By stationary
phase,
\[
\begin{aligned}
  \mathcal{A}(\infty)(\leq1)
  &=
  \bigoplus_{x\in\mathbb{A}^1}
  \mathcal{L}_{\psi_x}\otimes
  \operatorname{FTloc}(x,\infty)
  \bigl(\mathcal{F}(x)/\mathcal{F}(x)^{I(x)}\bigr),\\
  \mathcal{A}(\infty)(>1)
  &=
  \operatorname{FTloc}(\infty,\infty)(\mathcal{F}(\infty)(>1))
  =
  0,
\end{aligned}
\]
since \(\mathcal{F}(\infty)\) is tame. Since each \(\mathcal{F}(x)\) is tame,
we have
\[
  \operatorname{rank}
  \operatorname{FTloc}(x,\infty)
  \bigl(\mathcal{F}(x)/\mathcal{F}(x)^{I(x)}\bigr)
  =
  \operatorname{rank}\bigl(\mathcal{F}(x)/\mathcal{F}(x)^{I(x)}\bigr).
\]
Thus
\[
\begin{aligned}
  \operatorname{rank}M(\infty)
  &=
  \operatorname{rank}\mathcal{A}(0)
  =
  \operatorname{rank}\mathcal{A}
  =
  \operatorname{rank}\mathcal{A}(\infty)\\
  &=
  \sum_{x\in\mathbb{A}^1}
  \operatorname{rank}\bigl(\mathcal{F}(x)/\mathcal{F}(x)^{I(x)}\bigr)
  =
  \sum_{s\in\mathbb{A}^1-U}
  \operatorname{rank}\bigl(\mathcal{F}(s)/\mathcal{F}(s)^{I(s)}\bigr).
\end{aligned}
\]
The second formula for \(\operatorname{rank}M(\infty)\) is term by term equal
to this one: equate ranks in the isomorphism of \(I(s)\)-representations
\[
  \bigl(\mathcal{F}(s)/\mathcal{F}(s)^{I(s)}\bigr)
  \otimes\mathcal{L}_{\chi(x-s)}
  \cong
  \mathcal{G}(s)/\mathcal{G}(s)^{I(s)}.
\]
\end{proof}

\sourceheading[cor:3.3.7]{Corollary 3.3.7 (rank formula)}
Hypotheses and notations as in \hyperref[thm:3.3.3]{Theorem~3.3.3} above,
suppose in addition that \(K\) is tamely ramified everywhere. Then so is
\(K*_{\mathrm{mid}+}j_*\mathcal{L}_{\chi}[1]\), and their ranks and local
monodromies are related by the formulas
\[
\begin{aligned}
  \operatorname{rank}\mathcal{F}
  &=
  \sum_{s\in\mathbb{A}^1-U}
  \operatorname{rank}\bigl(\mathcal{G}(s)/\mathcal{G}(s)^{I(s)}\bigr)
  -
  \operatorname{rank}\bigl((\mathcal{G}(\infty)\otimes
  \mathcal{L}_{\overline{\chi}})^{I(\infty)}\bigr),\\
  \operatorname{rank}\mathcal{G}
  &=
  \sum_{s\in\mathbb{A}^1-U}
  \operatorname{rank}\bigl(\mathcal{F}(s)/\mathcal{F}(s)^{I(s)}\bigr)
  -
  \operatorname{rank}\bigl((\mathcal{F}(\infty)\otimes
  \mathcal{L}_{\chi})^{I(\infty)}\bigr).
\end{aligned}
\]

\begin{proof}
These follow immediately from the formulas
\[
\begin{aligned}
  \mathcal{F}(\infty)&=M(\infty)/M(\infty)^{I(\infty)},\\
  \mathcal{G}(\infty)&=
  M(\infty)\otimes\mathcal{L}_{\chi}/
  (M(\infty)\otimes\mathcal{L}_{\chi})^{I(\infty)},
\end{aligned}
\]
and
\[
\begin{aligned}
  \operatorname{rank}M(\infty)
  &=
  \sum_{s\in\mathbb{A}^1-U}
  \operatorname{rank}\bigl(\mathcal{F}(s)/\mathcal{F}(s)^{I(s)}\bigr),\\
  \operatorname{rank}M(\infty)
  &=
  \sum_{s\in\mathbb{A}^1-U}
  \operatorname{rank}\bigl(\mathcal{G}(s)/\mathcal{G}(s)^{I(s)}\bigr).
\end{aligned}
\]
\end{proof}

\section{Some open questions about local Fourier Transform}\label{sec:3.4}

\sourcepara{3.4.1}
We know that if \(M\) is an irreducible \(I(0)\)-representation with all slopes
\(a/b\), and dimension \(b\), then \(\operatorname{FTloc}(0,\infty)M\) is an
irreducible \(I(\infty)\)-representation with all slopes \(a/(a+b)\), and
dimension \(a+b\). We also know that \(\operatorname{FTloc}(0,\infty)\) is an
equivalence of categories
\[
  I(0)\text{-representations}
  \cong
  I(\infty)\text{-representations with all slopes }<1.
\]
Denote by \((\operatorname{FTloc}(0,\infty))^{-1}\) a quasi-inverse. Then for
every nontrivial Kummer sheaf \(\mathcal{L}_{\chi}\), we get an
autoequivalence of the category of \(I(0)\)-representations, defined by

\sourcepara{3.4.1.1}
\[
  M\longmapsto
  (\operatorname{FTloc}(0,\infty))^{-1}
  \bigl(\mathcal{L}_{\overline{\chi}}\otimes
  \operatorname{FTloc}(0,\infty)(M)\bigr).
\]
This autoequivalence preserves slopes and dimensions, and on tame \(M\) is
given (cf. \cite[7.4.4.3]{Ka-ESDE}) by
\(M\mapsto\mathcal{L}_{\chi}\otimes M\). What is it? Is there a simple formula
for it?

\sourcepara{3.4.2}
The most naive hope is that this autoequivalence be given by
\(M\mapsto\mathcal{L}_{\chi}\otimes M\) on all \(I(0)\)-representations, or
equivalently that
\[
  \operatorname{FTloc}(0,\infty)(M\otimes\mathcal{L}_{\chi})
  \cong
  \mathcal{L}_{\overline{\chi}}\otimes
  \operatorname{FTloc}(0,\infty)(M)
\]
as \(I(\infty)\)-representations. This hope is false. Here is a simple
sequence of counterexamples. For each integer \(n\geq1\), take for \(M_n\) the
\(I(0)\)-representation attached to
\(\operatorname{inv}^*(\operatorname{Kl}_n(\psi;\rho_1,\ldots,\rho_n))\),
where \(\operatorname{Kl}_n(\psi;\rho_1,\ldots,\rho_n)\) is any of the rank
\(n\) Kloosterman sheaves discussed in \cite[7.4.1]{Ka-GKM}. By
\cite[8.6.1]{Ka-GKM}, we have a geometric isomorphism
\[
  j^*\operatorname{FT}_{\psi}
  \bigl(j_!\operatorname{inv}^*
  (\operatorname{Kl}_n(\psi;\rho_1,\ldots,\rho_n))\bigr)
  \cong
  \operatorname{Kl}_{n+1}(\psi;\mathbf{1},\rho_1,\ldots,\rho_n).
\]
Because \(\operatorname{Kl}_{n+1}(\psi;\mathbf{1},\rho_1,\ldots,\rho_n)\) has
all its \(\infty\)-slopes \(<1\) (they are all equal to \(1/(n+1)\)),
stationary phase shows that
\[
  \operatorname{FTloc}(0,\infty)(M_n)
  =
  \operatorname{Kl}_{n+1}(\psi;\mathbf{1},\rho_1,\ldots,\rho_n)(\infty).
\]
If we replace \(M_n\) by \(M_n\otimes\mathcal{L}_{\chi}\), we are looking at
\[
  M_n\otimes\mathcal{L}_{\chi}
  :=
  \operatorname{inv}^*
  (\operatorname{Kl}_n(\psi;\overline{\chi}\rho_1,\ldots,
  \overline{\chi}\rho_n))
\]
as \(I(0)\)-representation, so by the above \(\operatorname{FT}\) formula we
have
\[
  \operatorname{FTloc}(0,\infty)(M_n\otimes\mathcal{L}_{\chi})
  =
  \operatorname{Kl}_{n+1}(\psi;\mathbf{1},\overline{\chi}\rho_1,\ldots,
  \overline{\chi}\rho_n)(\infty).
\]
We claim that
\[
  \operatorname{FTloc}(0,\infty)(M_n\otimes\mathcal{L}_{\chi})
  \not\cong
  \mathcal{L}_{\overline{\chi}}\otimes
  \operatorname{FTloc}(0,\infty)(M_n)
\]
as \(I(\infty)\)-representations. This amounts to the statement that, as
\(I(\infty)\)-representations,
\[
  \operatorname{Kl}_{n+1}(\psi;\mathbf{1},\overline{\chi}\rho_1,\ldots,
  \overline{\chi}\rho_n)
  \not\cong
  \mathcal{L}_{\overline{\chi}}\otimes
  \operatorname{Kl}_{n+1}(\psi;\mathbf{1},\rho_1,\ldots,\rho_n).
\]
Indeed, the two sides have nonisomorphic determinants. To see this, use the
geometric isomorphism (\cite[7.4.1]{Ka-GKM})
\[
  \det(\operatorname{Kl}_n(\psi;\rho_1,\ldots,\rho_n))
  \cong
  \mathcal{L}_{\rho_1\cdots\rho_n+1},
\]
valid for \(n\geq2\). The ratio of the two determinants is thus
\(\mathcal{L}_{\chi}\), which is nontrivial on \(I(\infty)\) so long as
\(\chi\) is nontrivial.

\sourcepara{3.4.3}
There is a somewhat unsatisfactory ``formula'' for the autoequivalence
\hyperref[par:3.4.1.1]{(3.4.1.1)}, in terms of middle convolution and the
canonical extension. Let \(M\) be an irreducible nontame representation of
\(I(0)\). By the theory of the canonical extension \cite{Ka-LG}, there exists
a lisse sheaf \(\mathcal{F}_M\) on \(\mathbb{G}_m\) which is tame at
\(\infty\), and whose \(I(0)\)-representation is \(M\). This sheaf
\(\mathcal{F}_M\) is certainly irreducible on \(\mathbb{G}_m\), since already
its \(I(0)\)-representation is irreducible. Denoting by
\(j:\mathbb{G}_m\to\mathbb{A}^1\) the inclusion, the object
\[
  K_M:=j_*\mathcal{F}_M[1]
\]
is a perverse irreducible on \(\mathbb{A}^1\) which (being wild at \(0\)) is
visibly of type \((2d)\), i.e., neither punctual nor an
\(\mathcal{L}_{\psi_\alpha}[1]\) nor a translate of a
\(j_*\mathcal{L}_{\chi}[1]\). Since \(K_M\) is tame at \(\infty\) and lisse on
\(\mathbb{G}_m\), Laumon's stationary phase tells us that
\(\operatorname{FT}K_M\) is lisse on \(\mathbb{G}_m\), say
\(j^*\operatorname{FT}K_M=\mathcal{A}[1]\) with \(\mathcal{A}\) a lisse sheaf
on \(\mathbb{G}_m\), and
\[
  \mathcal{A}(\infty)
  =
  \operatorname{FTloc}(0,\infty)(M)
  =
  \operatorname{FTloc}(0,\infty)(\mathcal{F}_M(0)).
\]
Now consider the object
\[
  L:=K_M*_{\mathrm{mid}+}j_*\mathcal{L}_{\chi}[1],
\]
whose \(\operatorname{FT}\) is also lisse on \(\mathbb{G}_m\), with
\(j^*\operatorname{FT}L=\mathcal{A}\otimes\mathcal{L}_{\overline{\chi}}[1]\).
The object \(L\) is lisse on \(\mathbb{G}_m\), and tame at \(\infty\) (by
\hyperref[thm:3.3.5]{Theorem~3.3.5}, \((2d)\)), since \(K_M\) is lisse on
\(\mathbb{G}_m\), and tame at \(\infty\). Say \(j^*L=\mathcal{N}[1]\), with
\(\mathcal{N}\) a lisse sheaf on \(\mathbb{G}_m\). Then by stationary phase,
\[
  (\mathcal{A}\otimes\mathcal{L}_{\overline{\chi}})(\infty)
  =
  \operatorname{FTloc}(0,\infty)(\mathcal{N}(0)).
\]
Thus the autoequivalence
\[
  M\longmapsto
  (\operatorname{FTloc}(0,\infty))^{-1}
  \bigl(\mathcal{L}_{\overline{\chi}}\otimes
  \operatorname{FTloc}(0,\infty)(M)\bigr)
\]
is given on nontame irreducibles \(M\) by
\[
  M\longmapsto\text{the }I(0)\text{-representation of }
  (j_*\mathcal{F}_M[1])*_{\mathrm{mid}+}(j_*\mathcal{L}_{\chi}[1]).
\]

\sourcepara{3.4.4}
We can also pose a similar question about \(\operatorname{FTloc}(\infty,\infty)\).
If \(M\) is an irreducible \(I(\infty)\)-representation with all slopes
\((a+b)/b>1\) and dimension \(b\), then \(\operatorname{FTloc}(\infty,\infty)(M)\)
is an irreducible \(I(\infty)\)-representation with all slopes
\((a+b)/a>1\) and dimension \(a\). We know that
\(\operatorname{FTloc}(\infty,\infty)\) is an autoequivalence of
\[
  \{I(\infty)\text{-representations with all slopes }>1\},
\]
with quasi-inverse
\([x\mapsto -x]^*\operatorname{FTloc}(\infty,\infty)\). For every nontrivial
Kummer sheaf \(\mathcal{L}_{\chi}\), we get another autoequivalence by

\sourcepara{3.4.4.1}
\[
  M\longmapsto
  [x\mapsto -x]^*\operatorname{FTloc}(\infty,\infty)
  \bigl(\mathcal{L}_{\overline{\chi}}\otimes
  \operatorname{FTloc}(\infty,\infty)(M)\bigr).
\]
This autoequivalence preserves slopes and dimensions. What is it? Is there a
simple formula for it?

\sourcepara{3.4.5}
The most naive hope is that the autoequivalence
\hyperref[par:3.4.4.1]{(3.4.4.1)} be given by
\(M\mapsto\mathcal{L}_{\chi}\otimes M\) on all \(I(\infty)\)-representations
with all slopes \(>1\), or equivalently that
\[
  \operatorname{FTloc}(\infty,\infty)(M\otimes\mathcal{L}_{\chi})
  \cong
  \mathcal{L}_{\chi}\otimes\operatorname{FTloc}(\infty,\infty)(M)
\]
as \(I(\infty)\)-representations. This hope is false. Here is a simple
sequence of counterexamples. For each \(n\geq3\) which is prime to \(p\), take
for \(M_n\) the \(I(\infty)\)-representation attached to
\(\mathcal{L}_{\psi(x^n)}\). For any Kummer sheaf \(\mathcal{L}_{\chi}\),
\(\operatorname{FTloc}(\infty,\infty)(M_n\otimes\mathcal{L}_{\chi})\) has rank
\(n-1\), and all slopes \(n/(n-1)\). We claim that

\sourcepara{3.4.5.1}
\[
  \det\bigl(\operatorname{FTloc}(\infty,\infty)
  (M_n\otimes\mathcal{L}_{\chi})\bigr)
  =
  \mathcal{L}_{\chi}
\]
as \(I(\infty)\)-representation. Admit this for a moment. Then
\(\operatorname{FTloc}(\infty,\infty)(M_n)\) has rank \(n-1\) and trivial
determinant for \(n\geq3\) prime to \(p\), so our naive hope would force
\(\operatorname{FTloc}(\infty,\infty)(M_n\otimes\mathcal{L}_{\chi})\) to have
determinant \(\mathcal{L}_{\chi}^{n-1}\) rather than \(\mathcal{L}_{\chi}\).
To show that \hyperref[par:3.4.5.1]{(3.4.5.1)} holds, consider first the case
of trivial \(\chi\). Then \(\operatorname{FT}_{\psi}(\mathcal{L}_{\psi(x^n)})\)
is lisse on \(\mathbb{A}^1\) of rank \(n-1\), and has all \(\infty\)-slopes
\(n/(n-1)\). As explained in \cite[9.2]{Ka-GKM}, its restriction to
\(\mathbb{G}_m\) descends through the \(n\)th power map, to a Kloosterman sheaf
of rank \(n-1\) on \(\mathbb{G}_m\). The determinant of this Kloosterman sheaf
is tame on \(\mathbb{G}_m\), because \(n-1\geq2\), and hence
\(\operatorname{FT}_{\psi}(\mathcal{L}_{\psi(x^n)})\) itself has a determinant
which is tame on \(\mathbb{G}_m\). But as
\(\operatorname{FT}_{\psi}(\mathcal{L}_{\psi(x^n)})\) is lisse on
\(\mathbb{A}^1\), so is its determinant, and hence
\(\det(\operatorname{FT}_{\psi}(\mathcal{L}_{\psi(x^n)}))=\mathbf{1}\) (being
both tame on \(\mathbb{G}_m\) and trivial at zero). As
\(\operatorname{FT}_{\psi}(\mathcal{L}_{\psi(x^n)})\) has all its
\(\infty\)-slopes \(>1\),
\[
  \operatorname{FTloc}(\infty,\infty)(M_n)
  =
  \operatorname{FT}_{\psi}(\mathcal{L}_{\psi(x^n)})(\infty)
\]
as \(I(\infty)\)-representation, and hence
\(\operatorname{FTloc}(\infty,\infty)(M_n)\) has trivial determinant, as
asserted.

Now consider the case when \(\chi\) is nontrivial. In this case,
\(\operatorname{FT}_{\psi}(\mathcal{L}_{\psi(x^n)}\otimes\mathcal{L}_{\chi})\)
is lisse on \(\mathbb{A}^1\) of rank \(n\), and its \(\infty\)-slopes are
\(0\) once and \(n/(n-1)\) repeated \(n-1\) times. As explained in
\cite[9.2.2]{Ka-GKM}, the restriction to \(\mathbb{G}_m\) of
\(\operatorname{FT}_{\psi}(\mathcal{L}_{\psi(x^n)}\otimes\mathcal{L}_{\chi})\)
descends through the \(n\)th power map to a hypergeometric sheaf of type
\((n,1)\). The determinant of such a hypergeometric sheaf is necessarily tame
on \(\mathbb{G}_m\) (because \(n\geq3\)), and hence
\(\operatorname{FT}_{\psi}(\mathcal{L}_{\psi(x^n)}\otimes\mathcal{L}_{\chi})\)
has trivial determinant (because this determinant is simultaneously lisse on
\(\mathbb{A}^1\) and tame on \(\mathbb{G}_m\)). By stationary phase, the
\(I(\infty)\)-representation of
\(\operatorname{FT}_{\psi}(\mathcal{L}_{\psi(x^n)}\otimes\mathcal{L}_{\chi})\)
is a direct sum
\[
  \operatorname{FT}_{\psi}(\mathcal{L}_{\psi(x^n)}\otimes\mathcal{L}_{\chi})(\infty)
  =
  \mathcal{L}_{\overline{\chi}}
  \oplus
  \operatorname{FTloc}(\infty,\infty)(M_n\otimes\mathcal{L}_{\chi}).
\]
Taking determinants, we find \hyperref[par:3.4.5.1]{(3.4.5.1)}, as required.

\sourcepara{3.4.6}
Just as in \hyperref[par:3.4.3]{(3.4.3)} above, we can give a somewhat
unsatisfactory ``formula'' for the autoequivalence
\hyperref[par:3.4.4.1]{(3.4.4.1)} in terms of canonical extensions. In terms
of the canonical extension of an \(I(\infty)\)-representation \(M\) with all
\(\infty\)-slopes \(>1\) to a lisse sheaf \(\mathcal{G}_M\) on
\(\mathbb{G}_m\) which is tame at zero, we get a description of the
autoequivalence \hyperref[par:3.4.4.1]{(3.4.4.1)} as
\[
  M\longmapsto
  \text{the slope }>1\text{ part of the }I(\infty)\text{-representation of }
  (j_*\mathcal{G}_M[1])*_{\mathrm{mid}+}(j_*\mathcal{L}_{\chi}[1]).
\]
